\section{本章小结}
\label{sec:opt-sto-summary}

章首的问题是：一百万行的最小二乘，若每步只读一行，单步方向可能指向上坡，许多步合起来是否仍在下降，又会停在哪里。本章的回答以有限步的期望误差为主线。只读一项的随机梯度只要在给定历史下条件无偏，代入光滑上界取条件期望后，期望下降与噪声上升就分成正比于$\eta$和$\eta^2$的两项。强凸时，常数步长使期望误差几何下降到一个与$\eta\sigma_\star^2/\mu$成正比的噪声邻域；递减步长$2/(\mu(t+t_0))$把它推向零，速率为$1/t$；凸目标的加权平均点和非凸目标的随机返回点分别有$1/\sqrt T$量级的界。这些界针对指定的返回对象，是期望意义下的结论；单条路径的几乎必然收敛属于随机逼近理论。

有限和、期望与变量分块三种结构分别支持三种节省。有限和允许只读一部分项，也允许反复访问同一项，因此可以用快照（SVRG）或逐项记忆（SAGA）把读数方差降到随迭代趋于零，在每项凸光滑、目标强凸时恢复线性收敛；期望目标只能不断抽取新样本，方差只能靠批量按$1/b$压低。非均匀抽样必须按抽样概率的倒数重加权才保持无偏，它可能改善步长上限，也可能增大噪声。历史读数还可以给出尺度：AdaGrad、RMSProp、Adam等逐坐标方法在问题的难点与坐标轴对齐时收益明显，旋转之后便失去优势；Adafactor压缩统计的存储，Shampoo与Muon在矩阵参数上处理行列方向，它们的成本和保证都必须连同各自的条件说明。变量分块的收益依赖于块常数小于整体光滑常数，失败来自耦合：耦合约束、不可分的非光滑项，以及分别凸而不联合凸的目标都可能让逐块精确的算法停在非最优点。

几个界限贯穿全章。条件偏差，无论来自裁剪、有偏抽样还是不精确的子问题，都在误差界中留下不随步长、批量和迭代次数消失的项；替代梯度连偏差都谈不上，只能回到实际前向任务上评价。算法的输出对单个样本替换的敏感程度可以在每项凸光滑的条件下用有限步界控制，但这不是泛化结论。最后，比较方法时要先固定目标的归一化、误差对象、返回点和停止准则，再统计读取、计算与存储的总量。长期过程的收敛归第~\ref{chap:stochastic-processes-and-approximation}章，数据误差与泛化由后续讨论学习理论的部分承担。
